\documentclass[a4paper,10pt]{article}

\usepackage[utf8]{inputenc}
\usepackage[T1]{fontenc}
\usepackage[portuguese]{babel}

\usepackage{fontspec}
\setmainfont{Times New Roman}

\usepackage{geometry}
\usepackage{parskip}
\usepackage{microtype}
\usepackage{enumitem}
\usepackage{titlesec}
\usepackage{array}
\usepackage{tabularx}

\usepackage{hyperref}
\usepackage{bookmark}
\usepackage{xurl}

\geometry{top=1.5cm, bottom=1.5cm, left=1.5cm, right=1.5cm} % Margens da página
\setcounter{secnumdepth}{0}                                 % Remove numeração de seção
\setlist[itemize]{
    leftmargin=0.75em, % Recuo da lista (mais perto da margem da página)
    itemsep=0.2em,     % Espaço entre itens
    topsep=0.25em,     % Espaço antes/depois da lista
    parsep=0em,        % Espaço entre parágrafos no mesmo item
    partopsep=0em      % Espaço extra ao iniciar a lista
}

\pagestyle{empty}

\hypersetup{
    pdftitle={CV Kevin Castro},      % Título do PDF
    pdfauthor={Kevin Castro},        % Autor do PDF
    colorlinks=true,          % Usa cor nos links (sem caixa)
    linkcolor=black,          % Cor de links internos
    urlcolor=black,           % Cor de URLs
    citecolor=black,          % Cor de citações
    bookmarksdepth=2          % Barra lateral: seção e subseção
}

\titleformat{\section}
{\Large\bfseries}
{}
{0em}
{}
[\titlerule\vspace{0.5ex}]
\titleformat{\subsection}
{\normalsize\bfseries}
{}
{0em}
{}
\titlespacing*{\subsection}{0pt}{0.35em}{0.15em}

\newcounter{cventry}
\newcommand{\cventry}[4]{%
    \refstepcounter{cventry}%
    \phantomsection%
    \pdfbookmark[2]{#1}{cventry-\thecventry}%
    \noindent\begin{tabularx}{\textwidth}{@{}>{\raggedright\arraybackslash}X >{\raggedleft\arraybackslash}X@{}}
    \textbf{#1} & #2 \\
    \textit{#3} & \textit{#4} \\
    \end{tabularx}
}

\begin{document}

\begin{center}
    {\LARGE \textbf{Kevin Castro}}
    \\ [0.1cm]
    {\normalsize Desenvolvedor Mobile}
    \\ [0.1cm]
    São Paulo, Brasil
    {\textbullet}
    \href{mailto:kevingcm.dev@gmail.com}{kevingcm.dev@gmail.com}
    {\textbullet}
    \href{https://linkedin.com/in/kevingcm}{linkedin.com/in/kevingcm}
    {\textbullet}
    \href{https://github.com/kevingcm}{github.com/kevingcm}
\end{center}

\section{Habilidades Técnicas}
    \begin{itemize}
        \item \textbf{Idiomas:} Português (Nativo), Inglês (C1)
        \item \textbf{Soft Skills:} Trabalho em equipe, Resolução de problemas, Atenção aos detalhes, Aprendizagem rápida, Gestão de tempo, Comunicação
        \item \textbf{Linguagens de Programação:} Dart, Python, JavaScript, HTML, CSS, Kotlin, Swift
        \item \textbf{Frameworks:} Flutter
        \item \textbf{Ferramentas:} Git, GitHub, PostgreSQL, Visual Studio Code, Figma
    \end{itemize}

\section{Educação}
    \cventry{Instituto Federal de Rondônia}{Porto Velho, Brasil}{Técnico em Informática}{2021 - 2023}
    \begin{itemize}
        \item \textbf{Disciplinas relevantes:} HTML \& CSS, Python, JavaScript, Git \& GitHub, Programação Orientada a Objetos, Banco de Dados Relacional \& SQL
    \end{itemize}

\section{Experiência}
    \cventry{Kmon! Climb Club, App de Escalada}{Projeto Pessoal}{Arquitetura full-stack \& produto}{Abr 2026 - Presente}
    \begin{itemize}
        \item Projetei e desenvolvi do zero um app completo com Flutter + Supabase, gerenciando mais de 50 migrações de banco de dados e definindo toda a arquitetura de dados.
        \item Implementei Row-Level Security (RLS) completo no PostgreSQL, políticas de acesso por tabela e controle de privacidade em três níveis.
        \item Construí um grafo social completo, com sistema de seguidores, feed de atividades, check-ins com fotos e curtidas e comentários.
        \item Desenvolvi busca com debounce de 250ms, filtros combinados e integração com Google Maps para exibir locais próximos.
    \end{itemize}

    \cventry{Instituto Federal de Rondônia}{Porto Velho, Brasil}{Monitoria de Programação Orientada a Objetos}{Mar 2023 - Out 2023}
    \begin{itemize}
        \item Ofereci suporte online diário aos alunos com conceitos e exercícios de programação em Python, todos os 7 dias da semana.
        \item Realizei sessões de reforço presenciais duas vezes por semana focadas na resolução de exercícios e desenvolvimento de lógica de programação.
        \item Apoiei o professor na revisão de atividades e trabalhos, contribuindo para a padronização e qualidade das avaliações.
    \end{itemize}

\end{document}
